\begin{figure*}[!htb]
  \begin{subfigure}[b]{0.49\textwidth}
    \centering
    \includegraphics[width=\textwidth]{figures/PR-10-10-0.2-1.png}
    \caption{1\% hits}
  \end{subfigure}
  \hfill
  \begin{subfigure}[b]{0.49\textwidth}
    \centering
    \includegraphics[width=\textwidth]{figures/PR-10-10-0.2-5.png}
    \caption{5\% hits}
  \end{subfigure}

  \begin{subfigure}[b]{0.49\textwidth}
    \centering
    \includegraphics[width=\textwidth]{figures/PR-10-10-0.2-10.png}
    \caption{10\% hits}
  \end{subfigure}
  \hfill
  \begin{subfigure}[b]{0.49\textwidth}
    \centering
    \includegraphics[width=\textwidth]{figures/PR-10-10-0.2-20.png}
    \caption{20\% hits}
  \end{subfigure}

  \begin{subfigure}[b]{0.49\textwidth}
    \centering
    \includegraphics[width=\textwidth]{figures/PR-10-10-0.2-30.png}
    \caption{30\% hits}
  \end{subfigure}
  \hfill
  \begin{subfigure}[b]{0.49\textwidth}
    \centering
    \includegraphics[width=\textwidth]{figures/PR-10-10-0.2-40.png}
    \caption{40\% hits}
  \end{subfigure}

  \caption{Comparison of PR curves for screening experiments with
    different hit rates (1\%, 5\%, 10\%, 20\%, 30\%, 40\%), using
    layouts with 10 positive and 10 negative controls in the presence
    of strong bowl-shaped plate effects. At low hit rates (1\%--5\%)
    PR-AUC is expectedly low due to class imbalance.}
    \label{fig:screening-pr-10-10}
\end{figure*}


%%% Local Variables: 
%%% mode: latex
%%% TeX-master: "../supplement.tex"
%%% End: 